\begin{abstract}
Flow matching has emerged as a powerful generative framework for visuomotor robot policies, offering stable training and high-quality multi-modal action generation. However, existing flow matching policies suffer from two orthogonal limitations: (1)~\emph{temporal blindness}---each decision is made independently without awareness of task phase or trajectory history; and (2)~\emph{spatial information loss}---global average pooling discards fine-grained spatial structure critical for precision tasks.

We propose the \textbf{\methodfull{} (\method{})} framework, which addresses both limitations through two self-supervised strategies. For temporal awareness, \method{} unifies three complementary perspectives: a \emph{Self-Predictive Latent Dynamics} model whose prediction error serves as an implicit phase-transition indicator (\emph{Future}); a \emph{Progress Predictor} with \emph{Dual FiLM} conditioning trained via ordinal ranking (\emph{Present}); and a lightweight \emph{Episodic Memory Bank} with same-episode context retrieval (\emph{Past}). For spatial awareness, \method{} introduces \textbf{\fasafull{} (\fasa{})}, a self-supervised module that discovers task-relevant regions by detecting inter-frame feature changes---decomposing spatial features into dynamic and static streams, encoding change centroids and directions, and extracting spatial keypoint coordinates---all without any external teacher model.

Both strategies require no manual annotations and integrate seamlessly into the VITA flow matching architecture. We evaluate \method{} on a challenging 6-DOF robotic S-curve seam tracking task, demonstrating significant improvements in tracking precision with complementary contributions from temporal and spatial awareness.
\end{abstract}
